%==============================================================================
% Research Methods - Group assignment - pitch
%==============================================================================

\documentclass[aspectratio=169]{beamer}

% You can increase the font size a bit with the option `14pt'

%==============================================================================
% Preamble
%==============================================================================

%---------- Layout ------------------------------------------------------------

\usetheme{hogent}

% Choose the background color theme, hgwhite (white background) or hgblack
% (black background)
\usecolortheme{hgwhite}

\setbeamercovered{transparent} % make invisible elements transparent

% override logo for specific slides
\newcommand{\whitelogo}{%
    \setbeamertemplate{footline}{%
        \begin{beamercolorbox}[wd=\paperwidth, dp=0.3cm]{footline}
            \usebeamercolor{page number in foot}
            \usebeamerfont{page number in foot}
            \hspace{1em}\insertframenumber/\inserttotalframenumber
            \hfill
            \begin{tikzpicture}[remember picture, overlay]
                \node[
                anchor=south east,
                shift={(-\the\wd\logopart, \the\wd\logopart)},
                inner sep=0,
                outer sep=0
                ] at (current page.south east) {%
                    
\includegraphics[
                    height=\the\ht\logopart
                    ]{hogent_logo_white}%
                };
            \end{tikzpicture}
        \end{beamercolorbox}
    }%
}

\definecolor{colorone}{HTML}{F0F3FA}
\definecolor{colortwo}{HTML}{D5DEEF}
\definecolor{colorthree}{HTML}{B1C9EF}
\definecolor{colorfour}{HTML}{8AAEE0}
\definecolor{colorfive}{HTML}{628ECB}
\definecolor{colorsix}{HTML}{395886}

% modify margin
%\setbeamersize{text margin left=3cm,text margin right=3cm}
%\setbeamertemplate{frametitle}{
%    \vspace{0.3cm}
%    \hspace{0.3cm}\insertframetitle
%}

%---------- Packages ----------------------------------------------------------
% You can add additional packages if you need them.

\usepackage[dutch]{babel}      % Dutch language support (remove or change to
                               % "english" for English)

\usepackage{booktabs}          % Better looking tables
\usepackage{multirow,multicol} % Merge table cells
\usepackage{eurosym}           % Euro symbool
\usepackage{lipsum}            % Filler text
\usepackage{graphicx}          % Required for inserting images
\usepackage{tcolorbox}         % Color boxes
\usepackage{xcolor}            % Custom text color
\usepackage{comment}           % comment-blocks


% DRY colorbox styles:
\tcbset{
    columntitle/.style={
        colback=colorthree,
        colframe=colorthree,
        coltext=colorsix,
        before skip=5pt,
        after skip=5pt,
        boxsep=.5pt,
        height=.7cm,
        boxrule=0pt,
        arc=2mm,
        fontupper=\small
    },
    columntext/.style={
        colback=colorone,
        colframe=colorone,
        coltext=colorsix,
        before skip=5pt,
        after skip=5pt,
        boxsep=0pt,
        boxrule=0pt,
        arc=2mm,
        left=4pt,
        right=4pt,
        top=4pt,
        bottom=4pt
    }
}
%---------- Metadata ----------------------------------------------------------

\title{\textcolor{white}{Pitch onderzoeksvoorstellen.}}
\subtitle{\textcolor{white}{Research Methods, 25-26, groep 85}}
\author{\textcolor{white}{Ewout Van Mullem \and Manu Coppens \and Peter Cypers \and Stef Simons}}
\date{\textcolor{white}{\today}}

%==============================================================================
% Presentation content
%==============================================================================

\begin{document}

{

    % An image cut out in the shape of a letter is a typical design element in
    % HOGENT's corporate identity. You can replace the image (in this case
    % "blurred-background-cellphone-coffee-842554.jpg") and the letter (in this
    % case "R") with your own choices. The image should be placed in the
    % `graphics` folder of your project and can be a jpeg or png.
    %\setbeamertemplate{background}[imgletter]
    %    {this-is-fine-art.png}{R}
    \whitelogo
    \usebackgroundtemplate{
\includegraphics[width=\paperwidth,height=\paperheight]{RM-bg-02.jpg}}

    \begin{frame}
        \maketitle
    \end{frame}
}

\begin{frame}{Onderzoeksvoorstellen}
    
    \vspace{-5pt}
     \begin{itemize}
        \item[]<1> \begin{tcolorbox}[colback=colortwo, colframe=colorsix, width=.8\textwidth]
            \textcolor{colorsix}{\textbf{Specialisatiedomein:} Ontwikkeling van een modulaire en contextbewuste Linux-desktopomgeving voor gebruikers die regelmatig wisselen tussen verschillende computer-activiteiten.}
        \end{tcolorbox}
        \item[]<2> \begin{tcolorbox}[colback=colortwo, colframe=colorsix, width=.8\textwidth]
            \textcolor{colorsix}{\textbf{Interprofessioneel:} Ontwerp en evaluatie van een digitale ecologische voetafdruk-tracker ter ondersteuning van duurzaam dagelijks gedrag bij studenten aan HOGENT: gebruik van gamification-technieken ter bevordering van gebruikersretentie.}
        \end{tcolorbox}
    \end{itemize}
    
    
\end{frame}


{
\usebackgroundtemplate{
\includegraphics[width=\paperwidth,height=\paperheight]{RM-bg-02.jpg}}
%\setbeamercolor{background canvas}{bg=black}
\whitelogo
\begin{frame}
    \begin{columns}
        \column{\textwidth} \huge \textbf{\textcolor{white}{Ontwikkeling van een modulaire en contextbewuste Linux-desktopomgeving voor gebruikers die regelmatig wisselen tussen verschillende computer-activiteiten.}}
    \end{columns}
\end{frame}
}

\begin{frame}[t]
    \begin{tcolorbox}[
        colback=colorfive,
        colframe=colorfive,
        coltext=white,
        before skip=0pt,
        after skip=0pt,
        boxrule=0pt,
        arc=2mm]
        \small \textbf{Ontwikkeling van een modulaire en contextbewuste Linux-desktopomgeving voor gebruikers die regelmatig wisselen tussen verschillende computer-activiteiten.}
    \end{tcolorbox}
    \begin{columns}[t,onlytextwidth] % align items top, max-width textwidth
% KOLOM 1
        \column{0.32\textwidth}
        \begin{tcolorbox}[columntitle]
            \textbf{Inleiding}
        \end{tcolorbox}
        % removed height, minimal explicit padding in 4 directions
        \begin{tcolorbox}[columntext]
            \tiny \lipsum[1][1-4]
        \end{tcolorbox}
        \begin{tcolorbox}[columntitle]
            \textbf{Probleemstelling}
        \end{tcolorbox}
        \begin{tcolorbox}[columntext]
            \tiny \lipsum[1][1-2]
        \end{tcolorbox}
        \begin{tcolorbox}[columntitle]
            \textbf{Onderzoeksvraag}
        \end{tcolorbox}
        % todo: deelvragen mee opnemen?
        \begin{tcolorbox}[columntext]
            \tiny \lipsum[1][1]
        \end{tcolorbox}
% KOLOM 2
        \column{0.32\textwidth} 
        \begin{tcolorbox}[columntitle]
            \textbf{Methode}
        \end{tcolorbox}
        \begin{tcolorbox}[columntext]
            \large $\ldots$ Methode Uitwerking $\ldots$
        \end{tcolorbox}
        \begin{tcolorbox}[columntitle, height=1.1cm]
            %note: dit past niet in KOL 3 want de HoGent logo neemt veel plek weg
            \textbf{Verwachtte \\ Resultaten}
        \end{tcolorbox}
        \begin{tcolorbox}[columntext]
            \tiny \lipsum[1][1-4]
        \end{tcolorbox}
% KOLOM 3
        \column{0.32\textwidth}
        \begin{tcolorbox}[columntext, colback=white, colframe=white, boxrule=0pt, boxsep=0pt, left=0pt, right=0pt, bottom=0pt]
            \tiny \textbf{\textcolor{colorsix}{
                    Hogent | Toegepaste Informatica | \\Research Methods Groep 85
            }}
        \end{tcolorbox}
        \begin{tcolorbox}[columntitle, height=1.1cm]
            \textbf{Besluit \& \\ Meerwaarde}
        \end{tcolorbox}
        \begin{tcolorbox}[columntext]
            \tiny \lipsum[1][1]
        \end{tcolorbox}
    \end{columns}   
\end{frame}

{
    \usebackgroundtemplate{
\includegraphics[width=\paperwidth,height=\paperheight]{RM-bg-02.jpg}}
    %\setbeamercolor{background canvas}{bg=black}
    \whitelogo
    \begin{frame}
        \begin{columns}
            \column{\textwidth} \huge \textbf{\textcolor{white}{Ontwerp en evaluatie van een digitale ecologische voetafdruk-tracker ter ondersteuning van duurzaam dagelijks gedrag bij studenten aan HOGENT: \\ \large Gebruik van gamification-technieken ter bevordering van gebruikersretentie.}}
        \end{columns}
    \end{frame}
}

%note: de titel is zo uitgebreid dat er minder plek over is om text te schrijven
\begin{frame}[t]
    \begin{tcolorbox}[
        colback=colorfive,
        colframe=colorfive,
        coltext=white,
        before skip=0pt,
        after skip=0pt,
        boxrule=0pt,
        arc=2mm]
        \small \textbf{Ontwerp en evaluatie van een digitale ecologische voetafdruk-tracker ter ondersteuning van duurzaam dagelijks gedrag bij studenten aan HOGENT: Gebruik van gamification-technieken ter bevordering van gebruikersretentie.}
    \end{tcolorbox}
    \begin{columns}[t,onlytextwidth] % align items top, max-width textwidth
        % KOLOM 1
        \column{0.32\textwidth}
        \begin{tcolorbox}[columntitle]
            \textbf{Inleiding}
        \end{tcolorbox}
        % removed height, minimal explicit padding in 4 directions
        \begin{tcolorbox}[columntext]
            \tiny \lipsum[1][1-2]
        \end{tcolorbox}
        \begin{tcolorbox}[columntitle]
            \textbf{Probleemstelling}
        \end{tcolorbox}
        \begin{tcolorbox}[columntext]
            \tiny \lipsum[1][1-2]
        \end{tcolorbox}
        \begin{tcolorbox}[columntitle]
            \textbf{Onderzoeksvraag}
        \end{tcolorbox}
        % todo: deelvragen mee opnemen?
        \begin{tcolorbox}[columntext]
            \tiny \lipsum[1][1]
        \end{tcolorbox}
        % KOLOM 2
        \column{0.32\textwidth} 
        \begin{tcolorbox}[columntitle]
            \textbf{Methode}
        \end{tcolorbox}
        \begin{tcolorbox}[columntext]
            \large $\ldots$ Methode Uitwerking $\ldots$
        \end{tcolorbox}
        \begin{tcolorbox}[columntitle, height=1.1cm]
            %note: dit past niet in KOL 3 want de HoGent logo neemt veel plek weg
            \textbf{Verwachtte \\ Resultaten}
        \end{tcolorbox}
        \begin{tcolorbox}[columntext]
            \tiny \lipsum[1][1-4]
        \end{tcolorbox}
        % KOLOM 3
        \column{0.32\textwidth}
        \begin{tcolorbox}[columntext, colback=white, colframe=white, boxrule=0pt, boxsep=0pt, left=0pt, right=0pt, bottom=0pt]
            \tiny \textbf{\textcolor{colorsix}{
                    Hogent | Toegepaste Informatica | \\Research Methods Groep 85
            }}
        \end{tcolorbox}
        \begin{tcolorbox}[columntitle, height=1.1cm]
            \textbf{Besluit \& \\ Meerwaarde}
        \end{tcolorbox}
        \begin{tcolorbox}[columntext]
            \tiny \lipsum[1][1]
        \end{tcolorbox}
    \end{columns}   
\end{frame}

\begin{frame}
    \frametitle{Conclusies.}
    
    \vspace{-5pt}
    \begin{itemize}
        \item[]<1> \begin{tcolorbox}[colback=colortwo, colframe=colorsix, width=.8\textwidth]
            \textcolor{colorsix}{\textbf{Specialisatiedomein:} \\ \lipsum[1][1]}
        \end{tcolorbox}
        \item[]<2> \begin{tcolorbox}[colback=colortwo, colframe=colorsix, width=.8\textwidth]
            \textcolor{colorsix}{\textbf{Interprofessioneel:} \\
                \lipsum[1][1]}
        \end{tcolorbox}
    \end{itemize}
\end{frame}

\begin{frame}
    \frametitle{Bronnen}
    
    {\huge \textbf{Einde Presentatie.}}
\end{frame}

% tcolorbox info lengths in pt, cm, mm
\begin{comment}
    \begin{tcolorbox}[
        colback=<colour>,        % Background/fill colour of the box
        colframe=<colour>,       % Colour of the border/frame
        coltext=<colour>,        % Colour of the text inside the box
        boxrule=<length>,        % Thickness of the border; 0pt = no visible border
        arc=<length>,            % Radius of the rounded corners; 0pt = square corners
        height=<length>,         % Fixed total height of the box
        width=<length>,
        boxsep=<length>,         % Internal padding between content and box edge
        ]
    \end{tcolorbox}
\end{comment}

\begin{frame}
    % empty slide
\end{frame}

% Presentatie WIP
\begin{frame}
    \frametitle{Onderzoeksvoorstellen 1}
    
    \begin{itemize}
        \item[] \textbf{Specialisatiedomein:} Ontwikkeling van een modulaire en contextbewuste Linux-desktopomgeving voor gebruikers die regelmatig wisselen tussen verschillende computer-activiteiten.
        \item[] \textbf{Interprofessioneel:} Ontwerp en evaluatie van een digitale ecologische voetafdruk-tracker ter ondersteuning van duurzaam dagelijks gedrag bij studenten aan HOGENT: gebruik van gamification-technieken ter bevordering van gebruikersretentie.
    \end{itemize}
\end{frame}

\begin{frame}{Onderzoeksvoorstellen 2}
    \begin{block}{Specialisatiedomein}
        Ontwikkeling van een modulaire en contextbewuste Linux-desktopomgeving voor gebruikers die regelmatig wisselen tussen verschillende computer-activiteiten.
    \end{block}
    
    \begin{exampleblock}{Interprofessioneel}
        Ontwerp en evaluatie van een digitale ecologische voetafdruk-tracker ter ondersteuning van duurzaam dagelijks gedrag bij studenten aan HOGENT: gebruik van gamification-technieken ter bevordering van gebruikersretentie.
    \end{exampleblock}
\end{frame}

\begin{frame}{Onderzoeksvoorstellen 3}
    \begin{tcolorbox}[colback=blue!30]
        \textbf{Specialisatiedomein:} Ontwikkeling van een modulaire en contextbewuste Linux-desktopomgeving voor gebruikers die regelmatig wisselen tussen verschillende computer-activiteiten
    \end{tcolorbox}
    \begin{tcolorbox}[colback=green!30]
        \textbf{Interprofessioneel:} Ontwerp en evaluatie van een digitale ecologische voetafdruk-tracker ter ondersteuning van duurzaam dagelijks gedrag bij studenten aan HOGENT: gebruik van gamification-technieken ter bevordering van gebruikersretentie.
    \end{tcolorbox}
\end{frame}

\begin{frame}{Onderzoeksvoorstellen 4}
    \colorbox{hgblue}{\textbf{Specialisatiedomein:} Ontwikkeling van een modulaire en contextbewuste Linux-desktopomgeving voor gebruikers die regelmatig wisselen tussen verschillende computer-activiteiten.} \\
    \colorbox{hglightgreen}{\textbf{Interprofessioneel:} Ontwerp en evaluatie van een digitale ecologische voetafdruk-tracker ter ondersteuning van duurzaam dagelijks gedrag bij studenten aan HOGENT: gebruik van gamification-technieken ter bevordering van gebruikersretentie.}
\end{frame}

% best op 2 aparte slides: titel + Probleemstelling
\begin{frame}
    {\huge \textbf{Ontwikkeling van een modulaire en contextbewuste Linux-desktopomgeving voor gebruikers die regelmatig wisselen tussen verschillende computer-activiteiten.}}
\end{frame}

\begin{frame}
    \begin{columns}
        \column{\textwidth} \huge \textbf{Ontwikkeling van een modulaire en contextbewuste Linux-desktopomgeving voor gebruikers die regelmatig wisselen tussen verschillende computer-activiteiten.}
    \end{columns}
\end{frame}

\section{Ontwikkeling van een modulaire en contextbewuste Linux-desktopomgeving voor gebruikers die regelmatig wisselen tussen verschillende computer-activiteiten}


% Sjabloon
\begin{frame}
    \huge \textbf{Sjabloon Hogent}
\end{frame}

\begin{frame}
    \frametitle{Inhoud.}

    \tableofcontents
\end{frame}

\section{Titel voorstel 1.}

\begin{frame}
    \frametitle{Inleiding.}

    Some text
\end{frame}

{
\setbeamertemplate{background}[imgletter]
    {blurred-background-cellphone-coffee-842554.jpg}{I}

\begin{frame}

    {\huge \textbf{Voorbeeld van tekst door een afbeelding}}

\end{frame}
}

\subsection{Subsectie 1.1.}

\begin{frame}
    \frametitle{Titel.}

    \begin{itemize}
        \item Lijn één
        \item Lijn twee
        \item Lijn drie
    \end{itemize}
\end{frame}

\subsection{Subsectie 1.2.}

\begin{frame}
    \frametitle{Twee kolommen.}

    \begin{columns}[c]
        \column{.5\textwidth}
        \begin{itemize}
            \item Lijn 1
            \item Lijn 2
            \item Lijn 3
        \end{itemize}

        \column{.5\textwidth}
        \begin{itemize}
            \item Lijn 1
            \item Lijn 2
            \item Lijn 3
        \end{itemize}
    \end{columns}
\end{frame}

\begin{frame}
    \frametitle{Kleuren}

    % These are examples of the colours defined in HOGENT's corporate identity.
    \textbf{\textcolor{hgdarkgreen}{tekst}
    \textcolor{hgpink}{tekst}
    \textcolor{hgochre}{tekst}
    \textcolor{hgorange}{tekst}
    \textcolor{hgpurple}{tekst}
    \textcolor{hgblue}{tekst}
    \textcolor{hglightgreen}{tekst}
    \textcolor{hgbrown}{tekst}
    \textcolor{hggrey}{tekst}
    \textcolor{hgyellow}{tekst}}

    \textbf{\colorbox{hgdarkgreen}{tekst}
    \colorbox{hgpink}{tekst}
    \colorbox{hgochre}{tekst}
    \colorbox{hgorange}{tekst}
    \colorbox{hgpurple}{tekst}
    \colorbox{hgblue}{tekst}
    \colorbox{hglightgreen}{tekst}
    \colorbox{hgbrown}{tekst}
    \colorbox{hggrey}{tekst}
    \colorbox{hgyellow}{tekst}}
\end{frame}

\section{Titel voorstel 2.}

\begin{frame}
\frametitle{Het tweede deel.}

\begin{table}
    \begin{tabular}{lll}
    \toprule
                    & \textbf{lorem ipsum} & \textbf{lorem ipsum} \\
    \midrule
    dolor sit amet & $3.255$              & $3.255$ \\
    \midrule
    lorem ipsum    & $425,43 (13,07\%)$   & $425,43 (13,07\%)$ \\
    \midrule
    dolor sit amet & $0,00$               & $1.273,31$ ($45\%$ -- gemiddeld)\\
    \midrule
    \textbf{{\small Verhouding kost}} & \textbf{$1,53$} & \textbf{$2,72$}\\
    \bottomrule
    \end{tabular}

    \label{tab:voorbeeld}
    \caption{Deze tabel heeft een bijschrift}
    \end{table}
\end{frame}

\subsection{Sectie 2.1.}

\begin{frame}
\frametitle{Afbeeldingen.}
    \begin{figure}
        \caption{Deze afbeelding heeft een bijschrift.}

        \includegraphics[height=.8\textheight]
            {beverage-coffee-computer-984536.jpg}
        % Bron: https://www.pexels.com/photo/hand-on-cup-of-coffee-984536/
        \label{img:voorbeeld}
    \end{figure}
\end{frame}

% If the HOGENT logo in the bottom right corner of the slide conflicts with
% the content, you can hide it by adding the option [plain] to the frame, as
% shown below.
\begin{frame}[plain]
    \includegraphics[width=\textwidth]{beverage-coffee-computer-984536.jpg}
\end{frame}

{
\setbeamertemplate{background}[imgletter]%
  {beverage-coffee-computer-984536.jpg}{T}

\begin{frame}

    {\huge \textbf{This is a famous quote that says something about
    my point I want to make.}}

    \bigskip

    \textbf{Name Quotemaker}

\end{frame}
}

\begin{frame}[fragile]
    \texttt{monogespatieerd lettertype met ligaturen -> <= >=}

    \begin{verbatim}
        a <- 1
        b <- 2
        a + b
    \end{verbatim}
\end{frame}

\begin{frame}{Wiskundige formule}
  \begin{eqnarray}
    X_{t} = \alpha x_{t} + (1-\alpha)(X_{t-1} + b_{t-1}) & 0 \leq \alpha \leq 1 \\
    b_{t} = \beta(X_{t}-X_{t-1}) + (1-\beta)b_{t-1} & 0 \leq \beta \leq 1 
    \label{eq:doubleSmoothing}
  \end{eqnarray}
\end{frame}

\end{document}
